%==============================================================================%
% BAB 1: PENDAHULUAN
%==============================================================================%
 
%-----------------------------------------------------------------------------%
\chapter{\babSatu}
%-----------------------------------------------------------------------------%
 
\section{Latar Belakang Masalah}
%-----------------------------------------------------------------------------%
 
Di era digitalisasi saat ini, transformasi teknologi tidak hanya terfokus pada
sektor teknologi informasi murni, tetapi juga telah merambah secara masif ke
sektor agrikultur dan perkebunan. Pengelolaan sumber daya manusia dan
administrasi operasional yang efisien menjadi pilar utama bagi perusahaan
agrikultur untuk mempertahankan produktivitas di lapangan. Kebutuhan akan sistem
informasi yang cepat, akurat, dan terpusat menjadi sangat krusial guna
menghindari kebocoran data dan manipulasi administratif yang merugikan perusahaan
\cite{almansouri2024hris}.
 
PT Pratama Nusantara Sakti adalah perusahaan yang bergerak di bidang agrikultur,
tepatnya perkebunan tebu dan produksi gula, dengan wilayah operasional yang cukup
luas di Ogan Komering Ilir, Sumatera Selatan, serta bekerja sama dengan sekitar
53 vendor atau kontraktor penyedia tenaga kerja kontrak yang tersebar di
berbagai provinsi asal. Mengingat jumlah tenaga kerja yang mencapai ribuan
orang, mengelola administrasi harian di lapangan menjadi tantangan yang tidak
sepele bagi perusahaan. Setiap hari, tim lapangan harus menangani proses
pendataan karyawan baru, penerbitan surat izin keluar, hingga pengajuan klaim
penggantian biaya transportasi.
 
Saat ini, tim di lapangan masih menggunakan aplikasi berbasis Microsoft Access
untuk mengurus registrasi pekerja, cetak surat izin keluar (\textit{permit}),
hingga pengajuan klaim transportasi (\textit{transport claim}). Masalahnya,
sistem ini sering mengalami \textit{error} karena tim IT di lapangan kurang
familiar dalam mengoperasikan basis data Microsoft Access yang bersifat lokal dan
tidak memiliki antarmuka yang intuitif. Akibatnya, banyak data tenaga kerja yang
tidak terinput dengan benar ke sistem, dan sering kali terjadi duplikasi data
seperti Nomor Induk Kependudukan (NIK) yang ganda pada satu basis data
\cite{ravikanth2024deduplication}.
 
Selain masalah teknis, sistem yang lama juga belum memiliki mekanisme validasi
yang kuat. Masih ditemukan celah kecurangan dari pihak vendor, seperti klaim
transport yang diajukan lebih dari satu kali, atau manipulasi jumlah tenaga kerja
demi keuntungan pribadi. Hal-hal seperti ini jelas merugikan perusahaan secara
finansial. Sebagai gambaran, pada proses migrasi data ditemukan sekitar 20
tenaga kerja dengan rekaman duplikat pada sistem lama; dengan asumsi setiap
duplikasi tersebut berpotensi memungkinkan pengajuan klaim transportasi ganda
dengan rata-rata nominal sekitar Rp2.000.000 per tenaga kerja, potensi kerugian
finansial dari temuan tersebut diperkirakan mencapai sekitar Rp40.000.000.
Mengingat sistem lama tidak memiliki mekanisme validasi otomatis untuk mencegah
duplikasi semacam ini, potensi kerugian tersebut berisiko berulang pada setiap
periode migrasi atau audit data berikutnya apabila tidak segera diatasi.
Penerapan validasi relasional yang ketat pada tingkat sistem dinilai penting
untuk menjaga integritas data transaksi sekaligus menutup celah kecurangan
tersebut \cite{zulkarnain2026validasi}.
Oleh karena itu, muncul kebutuhan mendesak untuk melakukan migrasi
dari sistem lokal yang rentan tersebut ke sebuah platform berbasis web yang lebih
rapi dan terpusat, yang kemudian dinamakan \textit{HR Safety Campus}
\cite{susanti2025hris}.
Dalam pengembangannya, platform \textit{HR Safety Campus} dibangun di atas
kerangka kerja Laravel yang menerapkan pola arsitektur
\textit{Model-View-Controller} (MVC). Pola ini memisahkan logika bisnis,
pengelolaan data, dan antarmuka secara tegas sehingga struktur kode menjadi
lebih rapi dan mudah dikembangkan dalam jangka panjang \cite{rahmawati2024mvc}.
Selain itu, pemilihan Laravel dibandingkan PHP \textit{native} turut
mempertimbangkan efisiensi dalam pengembangan layanan pertukaran data pada sisi
\textit{back-end} \cite{ariyanto2024restapi}.
 
Untuk mengatasi celah manipulasi tersebut, sistem baru ini memerlukan mekanisme
pembatasan hak akses yang ketat berdasarkan peran pengguna di lapangan
\cite{mary2025rbac, yuricha2023rbac}. Selain itu, karena kondisi lapangan yang dinamis dan
perangkat yang bervariasi, antarmuka sistem harus dirancang agar sangat responsif
dan ringan saat diakses oleh tim IT di berbagai lokasi \cite{nandan2024tailwind}.
Walaupun namanya mengandung unsur \textit{safety}, fokus utama aplikasi ini
adalah mendigitalisasi registrasi tenaga kerja, mengotomatiskan cetak surat izin,
dan memperketat validasi klaim transportasi agar tidak ada lagi celah kecurangan
operasional.
 
%-----------------------------------------------------------------------------%
\section{Maksud dan Tujuan Kerja Magang}
%-----------------------------------------------------------------------------%
 
\subsection{Maksud Kerja}
 
\medskip
 
Pelaksanaan kerja magang di PT Pratama Nusantara Sakti sebagai
\textit{Full-Stack Developer} dimaksudkan untuk memberikan pengalaman nyata
mengenai siklus pengembangan perangkat lunak (\textit{Software Development Life
Cycle}) dalam lingkup skala \textit{enterprise}. Melalui kegiatan ini, kegiatan magang 
dimaksudkan untuk mempelajari secara langsung bagaimana menangani migrasi data dalam 
jumlah besar, merancang arsitektur keamanan akses berbasis peran, serta beradaptasi dengan 
ritme kerja profesional dan pola komunikasi langsung dengan klien.
 
\subsection{Tujuan Kerja}
 
\medskip
 
Selama menjalani program magang, terdapat sejumlah tujuan operasional yang
menjadi fokus penyelesaian masalah, antara lain:
 
\vspace{0.5em}
 
\begin{enumerate}
    \item Merancang struktur navigasi sistem (\textit{site map}) dan alur logika bisnis (\textit{flowchart}) untuk keenam fitur utama sistem \textit{HR Safety Campus} sebelum tahap implementasi, guna memastikan rancangan sistem terdefinisi dengan jelas sebelum proses pengkodean dimulai.
    \item Membangun antarmuka (\textit{frontend}) sistem registrasi tenaga kerja berbasis web menggunakan pendekatan \textit{utility-first} \cite{nandan2024tailwind} untuk menggantikan antarmuka Microsoft Access yang lama, dengan siklus pengembangan mengikuti metode \textit{Waterfall} \cite{ardiyanto2023waterfall}.
    \item Merancang dan mengimplementasikan arsitektur \textit{Role-Based Access Control} (RBAC) pada logika \textit{backend} untuk membatasi akses menu pencetakan surat izin (\textit{permit}) dan klaim transportasi sesuai dengan peran masing-masing pengguna \cite{mary2025rbac}.
    \item Melakukan pembersihan data (\textit{data cleansing}) dan migrasi
    sekitar 4.000 data tenaga kerja dari sistem lama ke basis data relasional
    MySQL, termasuk penanganan duplikasi NIK yang ditemukan selama proses
    migrasi berlangsung \cite{ravikanth2024deduplication}.
    \item Memperketat validasi data melalui sistem untuk meminimalisasi kesalahan
    input (\textit{human error}) dan menutup celah praktik kecurangan klaim
    oleh pihak vendor.
\end{enumerate}
 
%-----------------------------------------------------------------------------%
\section{Waktu dan Prosedur Kerja Magang}
%-----------------------------------------------------------------------------%
 
\subsection{Waktu Pelaksanaan Kerja}
 
\medskip
 
Kegiatan kerja magang dilaksanakan di PT Pratama Nusantara Sakti. Waktu
pelaksanaan dan sistem kerja diuraikan sebagai berikut:
 
\vspace{0.5em}
 
\begin{itemize}
    \item \textbf{Periode Magang:} 05 Januari 2026 sampai dengan 20 Mei 2026.
    \item \textbf{Posisi:} \textit{Full-Stack Developer}.
    \item \textbf{Jam Operasional dan Sistem Kerja:} Bekerja secara penuh dari
    kantor (\textit{Work From Office} / WFO) setiap hari kerja mulai pukul
    08.00 hingga 17.00 WIB. Seluruh kegiatan operasional harian, termasuk
    rapat dan koordinasi tim, dilaksanakan secara tatap muka di lingkungan
    kantor perusahaan.
\end{itemize}
 
\subsection{Prosedur Pelaksanaan Kerja}
 
\medskip
 
Prosedur pelaksanaan magang dibagi menjadi tiga tahapan utama:
 
\vspace{0.5em}
 
\begin{enumerate}
    \item \textbf{Tahap Pra-Magang:} Proses rekrutmen diawali dengan pengiriman
    dokumen lamaran dan portofolio melalui surat elektronik (\textit{e-mail})
    langsung kepada pihak \textit{Human Resources Department} (HRD)
    perusahaan. Setelah lolos seleksi awal, prosedur dilanjutkan dengan tahap
    wawancara \textit{user}. Evaluasi pada tahap ini tidak melibatkan
    \textit{live coding}, melainkan difokuskan pada pembedahan portofolio,
    rekam jejak penyelesaian proyek sebelumnya, serta pendalaman keahlian
    teknis secara lisan.
    \item \textbf{Tahap Pelaksanaan Magang:} Diawali dengan observasi alur kerja
    Microsoft Access di lapangan dan penyiapan lingkungan pengembangan lokal
    pada laptop. Pekerjaan kemudian dilanjutkan dengan eksekusi penulisan kode,
    pembersihan ribuan baris data CSV, perancangan fungsi \textit{backend}
    menggunakan kerangka kerja Laravel, hingga konfirmasi kelayakan data dengan
    pihak klien.
    \item \textbf{Tahap Pasca-Magang:} Tahap akhir difokuskan pada penyusunan
    dokumen laporan kerja magang sebagai bentuk pertanggungjawaban akademis
    kepada pihak universitas.
\end{enumerate}